\begin{proposition}[5.10.14]
\label{IV.2.s5.10.14-fr}
Sean $X$ un preesquema localmente noetheriano, $\mathcal F$ un
$\mathcal O_X$-módulo coherente de soporte igual a $X$.
\begin{enumerate}
\item[\textup{(i)}] Para que $\mathcal F$ verifique $(S_1)$, es necesario y suficiente
que sea $Z^{(1)}$-puro.
\item[\textup{(ii)}] Para que $\mathcal F$ verifique $(S_2)$, es necesario y suficiente
que sea $Z^{(2)}$-cerrado y $Z^{(1)}$-puro, o también que sea $Z^{(2)}$-cerrado
y carezca de ciclo primo asociado de codimensión 1.
\end{enumerate}
\end{proposition}

(i) Decir que $\mathcal F$ posee la propiedad $(S_1)$ significa que
$\mathcal F$ carece de ciclo primo asociado inmerso (5.7.5), o también que,
para todo $x\in\operatorname{Ass}(\mathcal F)$, se tiene
$\dim(\mathcal F_x)=0$ (3.1.4), es decir (5.1.12.1)
$\dim(\mathcal O_x)=0$; pero esto equivale a decir que
$\operatorname{Ass}(\mathcal F)$ no corta $Z^{(1)}$, y la conclusión se sigue
de (5.10.2).

(ii) Decir que $\mathcal F$ es $Z^{(2)}$-cerrado significa que
$\operatorname{prof}_{Z^{(2)}}(\mathcal F)\geq2$, o también que, para todo
$x\in X$, la relación $\dim(\mathcal F_x)\geq2$ implica
$\operatorname{prof}(\mathcal F_x)\geq2$; esto muestra que la propiedad
$(S_2)$ implica que $\mathcal F$ es $Z^{(2)}$-cerrado; implica además que
$\mathcal F$ verifica $(S_1)$, y por tanto carece de ciclo primo asociado
inmerso (5.7.5), y como $\operatorname{Supp}(\mathcal F)=X$, esto significa
también que todos los ciclos primos asociados a $\mathcal F$ tienen
codimensión 0. Recíprocamente, supongamos que $\mathcal F$ sea
$Z^{(2)}$-cerrado y carezca de ciclo primo asociado de codimensión 1; para ver
que $\mathcal F$ verifica $(S_2)$, queda por demostrar que, si $x\in X$ es tal
que $\dim(\mathcal F_x)=1$ (o, lo que viene a ser lo mismo,
$\dim(\mathcal O_x)=1$), entonces se tiene
$\operatorname{prof}(\mathcal F_x)=1$; pero, por hipótesis, la relación
$\dim(\mathcal O_x)=1$ implica $x\notin\operatorname{Ass}(\mathcal F)$, y
esta última relación equivale a $\operatorname{prof}(\mathcal F_x)\neq0$, es
decir, en este caso, a $\operatorname{prof}(\mathcal F_x)=1$. Si $\mathcal F$
es $Z^{(1)}$-puro, y por tanto verifica $(S_1)$, se ha observado más arriba que,
en virtud de la relación $\operatorname{Supp}(\mathcal F)=X$, todos los ciclos
primos asociados a $\mathcal F$ tienen codimensión 0, por lo que se aplica lo
anterior.

Se observará que puede ocurrir que $\mathcal F$ sea $Z^{(2)}$-cerrado y no
verifique $(S_1)$: así sucede, por ejemplo, cuando $X$ es de dimensión 1
(pues entonces $Z^{(2)}=\varnothing$, y todo $\mathcal O_X$-módulo es
$Z^{(2)}$-cerrado) y posee ciclos primos asociados inmersos.

Recordemos que en el capítulo III, en el estudio de la cohomología local, se
da una caracterización cohomológica de la propiedad $(S_n)$ para todo
$n\geq1$, que generaliza (5.10.14).

\begin{corollary}[5.10.15]
\label{IV.2.s5.10.15-fr}
Sean $X$ un preesquema localmente noetheriano, $\mathcal F$ un
$\mathcal O_X$-módulo coherente de soporte $X$. Se supone que $\mathcal F$
carece de ciclos primos asociados de codimensión 1 y que
$\mathcal F'=\mathcal H^0_{X/Z^{(2)}}(\mathcal F)$ sea coherente. Entonces:
\begin{enumerate}
\item[\textup{(i)}] $\mathcal F'$ verifica la propiedad $(S_2)$.

\oldpage[IV-2]{120}

\item[\textup{(ii)}] El conjunto $U$ de los $x\in X$ tales que $\mathcal F$ verifica
$(S_2)$ en el punto $x$ \textup{(5.7.2)} es abierto en $X$ y se tiene
$\operatorname{codim}(X\setmin U,X)\geq2$.
\end{enumerate}
\end{corollary}

(i) Se sabe (5.9.11) que $\mathcal F'$ es $Z^{(2)}$-cerrado, y además
$\operatorname{Supp}(\mathcal F')=X$, pues los puntos maximales de $X$
pertenecen a $X\setmin Z^{(2)}$ y en esos puntos
$\mathcal F'_x=\mathcal F_x\neq0$; por tanto, el soporte de $\mathcal F'$ es
denso en $X$, y como $\mathcal F'$ es coherente,
$\operatorname{Supp}(\mathcal F')$ es cerrado y, por tanto, igual a $X$. Queda
por ver que $\mathcal F'$ carece de ciclos primos asociados de codimensión 1.
Pero si $x\in\operatorname{Ass}(\mathcal F')$ y
$\dim(\mathcal F'_x)=\dim(\mathcal O_x)=1$, se tiene $x\in X\setmin Z^{(2)}$, de
modo que de la igualdad $\mathcal F'_x=\mathcal F_x$ se seguiría que
$x\in\operatorname{Ass}(\mathcal F)$, en contra de la hipótesis, lo que
termina de demostrar (i).

(ii) Se tiene $Z^{(n)}(X_x)=Z^{(n)}(X)\cap X_x$, con las notaciones de
(5.9.6), teniendo en cuenta (\textbf I, 2.4.2); por otra parte, la hipótesis de que
$\mathcal F$ carece de ciclos primos asociados de codimensión 1 implica la
misma hipótesis para $\widetilde{\mathcal F}_x$; para que
$\widetilde{\mathcal F}_x$ verifique $(S_2)$, es necesario y suficiente, por
tanto, en virtud de (5.10.14), que $\widetilde{\mathcal F}_x$ sea
$Z^{(2)}(X_x)$-cerrado; la afirmación (ii) se deduce, por tanto, de (5.10.6) y
de (5.10.10, (iii)).
